\section*{Limitations}

In this paper, we conduct a thorough research into the relationship between legal statutes and precedents. Specifically, we have made the first attempt (to the best of our knowledge) to solve the tasks of LSR and PCR simultaneously from the same query. All the prior works have either taken isolated approaches to solve the two tasks, or considered statute semantics while understanding PCR, but no work has tried to effectively solve both tasks by utilizing the inherent relationship between them. 

\noindent Our experiments have revealed that this could be a difficult exercise, since different types of features (lexical vs. semantic) are important for the two different tasks. The multi-task results are counter-intuitive, since despite the inherent relationship between statutes and precedents cited from the same query, independently trained models fare better in most settings. 
%Despite our best engineering efforts, such as, using variable learning rates for different layers, advanced negative sampling techniques including in-batch sampling, etc., we are not able to improve multi-task results for any of the models, even in the ensemble setup. Although we have discussed some of the possible reasons behind this in Section~\ref{sec:expt}, this needs more investigation and thorough studies.
However, we do find that a sequential/pipeline approach to solving the tasks can be beneficial for both. We observe this effect during transfer learning at the ensemble (retriever) phase, as well as two-stage prompting for the LLM (re-ranker) phase. We have observed these effects only with a specific set of models, i.e., Para-GNN (summaries) + BM25 as the ensemble retriever, and GPT-4.1 as the LLM re-ranker. More experiments with other models and investigations are needed to study the effects of such techniques across models and approaches.

%Another aspect of legal retrieval tasks, that need to be considered is the length and nature of the query input. As we mentioned in Section~\ref{sec:dataset/formulation}, although most prior works have taken the full query document text as the input, some works have considered only the facts of the case, or additionally, the arguments by the lawyers. This is indeed much closer to the practical setting -- in real life, for new legal situations, lawyers and judges will mostly have just the facts to identify statutes and precedents, and not full case documents~\citep{feng2024pcrsurvey}. In fact, most prior LSR approaches~\citep{chalkidis2019neural,paul2024legal} have worked only with the facts of the case. Indeed, there could be a significant difference in performance when considering just the facts vs. the full case document, since a lot of information that helps in PCR matching is actually missing from the facts. The amount of case information actually needed for performing LSR and PCR effectively needs to be studied more carefully in future work.

\noindent Certain design decisions of the dataset also merit discussion. Our candidate pools consist of many \textit{prominent} statutes and judgments, i.e., candidates that are frequently cited. This is a standard practice, according to prior works~\citep{paul2022ilsi,joshi2023ucreat}. This strategy ensures that we are able to cover a large number of queries, which is essential for both training and development. Although this leads to an inherent bias towards the frequent candidates, we ensure some diversity by considering a lot of candidates that have < 5 cited queries, and even some candidates that have no citations (held-out candidates). 

\noindent Additionally, as discussed in \S\ref{sec:dataset/overview}, the concept of relevance in the legal domain can be quite narrow, as in, all prior cases similar to the query are not necessarily cited. Similarly, only a particular statute from a family of similar statutes is usually applied based on the exact circumstances of the case. Utilizing just the cited candidates as the gold-standard can thus be restrictive, although this has been the standard practice~\citep{paul2022ilsi,joshi2023ucreat,li2024coliee2024}. To address this concern, we conducted annotation exercises with domain experts (see \S~\ref{sec:human-annotation}). We observed that the annotators not only differed from the cited gold-standard, but also among themselves, highlighting the subjectivity of relevance in legal retrieval. However, we observed that most of the model trends remain similar when considering just the cited candidates vs. union of citations and human annotations as the gold-standard, suggesting that the set of cited candidates, overall, can serve as an efficient proxy for the set of all possible relevant candidates.


\section*{Ethical Considerations}

\noindent In this work, we propose a system that allows for the retrieval of both statutes and precedents given a query case. Both these tasks are extremely crucial for the legal domain, and legal professionals regularly require technological assistance to reduce the search space of candidate statutes/precedents. These methods are designed to only provide relevant recommendations to assist the legal professionals, and are \textit{not} expected to be integrated directly into the decision-making process of the judicial system. 
Further, we ensured that all case documents used in our dataset \ilpcsr{} are publicly available. We also took steps to pre-process the documents by removing entity mentions that can lead to biases in the models.

